% Rozpis: 1_rocnik/tematicke_plany/Rozpis_vyuky_MAT_1_rocnik.tex, hodina 8;
% SPS.pdf, oddíl 1.2. Navazuje na číselnou osu 07; operace s intervaly až v 18.
% Učebnice: ucebnice/01_Zakladni_poznatky/Ciselne_obory/Číselné obory.tex,
% oddíl Intervaly. Značení s lomenými závorkami zachováno.
% Realisticky.cz: externi_zdroje/realisticky/01_Zakladni_poznatky/
% 03_Mnoziny/010308_Intervaly/010308_Lekce.pdf, úvod, Př. 1–3 a přehled intervalů.
% Metodická inspirace: rozdíl mezi spojitým intervalem a výběrem celých čísel;
% slovní podmínka -> osa -> symboly. Žádné množinové operace nepřebírány.
% Výklad, obrázky, konkrétní úlohy a řešení jsou vlastní tvorba.
\documentclass[11pt,a4paper]{article}
% Spolecny styl pro ucitelske pripravy MAT (Work/MAT)
% Verze 1.1, 2026-10-23
\NeedsTeXFormat{LaTeX2e}
\ProvidesPackage{priprava}[2026/10/23 v1.1 Sablona ucitelskych priprav MAT]

% Zakladni typografie a cestina
\RequirePackage[T1]{fontenc}
\RequirePackage[utf8]{inputenc}
\RequirePackage[czech]{babel}
\RequirePackage{lmodern}
\RequirePackage{microtype}
\RequirePackage{amsmath,amssymb}
\RequirePackage{xcolor}
\RequirePackage{enumitem}
\RequirePackage[most]{tcolorbox}
\RequirePackage{tabularx,booktabs,array}
\RequirePackage{fancyhdr}
\RequirePackage{needspace}
\RequirePackage{geometry}
\RequirePackage[hidelinks]{hyperref}

\geometry{a4paper,top=17mm,bottom=18mm,left=23mm,right=23mm}
\definecolor{notegray}{gray}{0.39}
\setlength{\parindent}{0pt}
\setlength{\parskip}{2.5pt plus 1pt minus 0.5pt}
\setlist[itemize]{topsep=2pt,itemsep=1pt,parsep=0pt,leftmargin=1.45em}
\setlist[enumerate]{topsep=2pt,itemsep=2pt,parsep=0pt,leftmargin=1.65em}
\setcounter{secnumdepth}{0}
\renewcommand{\arraystretch}{1.12}

% Zahlavi: vlevo rocnik | globalni tema, vpravo cislo: nazev hodiny
% Zapati zustava zcela prazdne.
\pagestyle{fancy}
\fancyhf{}
\renewcommand{\headrulewidth}{0.35pt}
\renewcommand{\footrulewidth}{0pt}
\setlength{\headheight}{14pt}
\newcommand{\nastavPripravu}[4]{%
  \fancyhead[L]{\small #1\enspace\textbar\enspace #2}%
  \fancyhead[R]{\small #3: #4}%
}

% Minimalisticke cernobile boxy. Kazda definice zvlast.
\tcbset{commonlesson/.style={enhanced,breakable=false,colback=white,colframe=black,
  colbacktitle=white,coltitle=black,boxrule=0.45pt,
  arc=0.8mm,outer arc=0.8mm,left=3mm,right=3mm,top=1.6mm,bottom=1.6mm,
  toptitle=1mm,bottomtitle=0.7mm,fonttitle=\bfseries,
  attach boxed title to top left={xshift=2.5mm,yshift=-1.15mm},
  boxed title style={colback=white,colframe=white,boxrule=0pt,left=1mm,right=1mm}}}
\newtcolorbox{zadaniBox}[1][]{commonlesson,title={Zadání},#1}
\newtcolorbox{prikladBox}[1][]{commonlesson,title={Příklad},#1}
\newtcolorbox{definiceBox}[1][]{commonlesson,title={Definice},#1}

% Spolecny box pro zadani a jeho reseni, bez mezery mezi castmi.
% Pouziti:
%   \begin{prikladReseniBox}
%     Zadani ulohy.
%     \tcblower
%     \begin{align*} ... \end{align*}
%   \end{prikladReseniBox}
% Spodni cast ma jemny sedy podklad a delici linku; lze v ni sazet
% vice radku matematiky. Obe casti maji jen jeden spolecny ramecek.
\tcbset{solutionlesson/.style={
  commonlesson,
  enhanced jigsaw,
  breakable,
  colbacklower=black!3,
  segmentation style={solid,draw=black!25,line width=0.35pt},
  fontlower=\small\color{notegray},
  before lower={\noindent\textbf{Řešení:}\enspace},
  leftlower=3mm,rightlower=3mm
}}
\newtcolorbox{prikladReseniBox}[1][]{solutionlesson,title={Příklad},#1}
\newtcolorbox{zadaniReseniBox}[1][]{solutionlesson,title={Zadání},#1}

% Pomucky pro autora pripravy: sedive poznamky; \reseni pro starsi pripravy
\newcommand{\poznamka}[1]{\par\begingroup\small\color{notegray}\itshape
  \textnormal{\textbf{Poznámka pro učitele.}}\enspace #1\par\endgroup}
\newcommand{\reseni}[1]{\par\begingroup\small\color{notegray}
  \textbf{Očekávané řešení:}\enspace #1\par\endgroup}
\newcommand{\blok}[1]{\needspace{5\baselineskip}\subsection*{#1}}
\newcommand{\vyrok}[1]{\uv{#1}}

\nastavPripravu{1. ročník}{Číselné obory}{08}{Intervaly}
\begin{document}

\poznamka{Orientačně: motivace a osa 4 min, význam závorek 9 min,
řešené převody 9 min, samostatná práce s kontrolou krajních bodů 11 min,
shrnutí 2 min; základ asi 35 min. Rezerva 10–15 min k výběru,
domácí práce 15–20 min. Každý žák kreslí vlastní osy, nesleduje jen hotové obrázky.}

\blok{Jak zapsat všechny přípustné hodnoty?}
Součástka smí mít délku alespoň $9{,}8$ mm a nejvýše $10{,}2$ mm.
Vyhovují krajní délky i všechna reálná čísla mezi nimi, nejen desetinná
čísla s jedním místem. Nekonečný výčet nahradíme zápisem intervalu.

\begin{definiceBox}
\textbf{Interval} je souvislá část číselné osy: spolu s libovolnými dvěma
svými čísly obsahuje všechna reálná čísla mezi nimi.
Zahrnutí krajních bodů vyznačujeme závorkami.
\end{definiceBox}

V následujících definicích předpokládáme $a,b\in\mathbb R$ a $a<b$.
V zápisu oddělujeme meze středníkem, aby se nepletl s desetinnou čárkou.

\begin{definiceBox}
\textbf{Uzavřený interval} $\langle a;b\rangle$ obsahuje právě všechna
reálná $x$, pro která $a\le x\le b$. Oba krajní body do něj patří.
\end{definiceBox}

\begin{definiceBox}
\textbf{Otevřený interval} $(a;b)$ obsahuje právě všechna reálná $x$,
pro která $a<x<b$. Ani jeden krajní bod do něj nepatří.
\end{definiceBox}

\begin{definiceBox}
\textbf{Polouzavřený interval} zahrnuje právě jeden krajní bod:
$\langle a;b)$ znamená $a\le x<b$, zatímco $(a;b\rangle$ znamená $a<x\le b$.
\end{definiceBox}

Na ose značíme zahrnutý kraj plným kolečkem a nezahrnutý prázdným.
Úsek mezi nimi zvýrazníme spojitou čarou, protože zahrnuje i necelá čísla.

\begin{prikladReseniBox}[title={P1: nejprve rozhodni o krajích},breakable=false]
Zapiš intervalem a nakresli všechna reálná čísla $x$ s vlastností $-2\le x<3$.
Rozhodni, zda vyhovují čísla $-2$, $0{,}5$ a $3$.
\tcblower
Jde o $\langle-2;3)$. První dvě čísla vyhovují, $3$ nikoli.
\begin{center}
\begin{tikzpicture}[x=1.3cm,y=0.7cm]
\draw[->] (-3.1,0)--(4,0);
\foreach \x in {-3,-2,-1,0,1,2,3}{\draw (\x,-0.08)--(\x,0.08) node[below=4pt] {$\x$};}
\draw[line width=1.4pt] (-2,0)--(3,0);
\fill (-2,0) circle (2.8pt);
\draw[fill=white,line width=0.9pt] (3,0) circle (2.8pt);
\end{tikzpicture}
\end{center}
Dosazení do původní podmínky je kontrolou závorek:
$-2\le-2<3$ platí, ale $-2\le3<3$ neplatí.
\end{prikladReseniBox}

\newpage
\blok{Když je omezení jen na jedné straně}
\begin{definiceBox}
\textbf{Neomezený interval} pokračuje alespoň jedním směrem bez konce.
Například $(a;+\infty)$ znamená $x>a$, $\langle a;+\infty)$ znamená $x\ge a$,
$(-\infty;b)$ znamená $x<b$ a $(-\infty;b\rangle$ znamená $x\le b$.
Vždy uvažujeme reálná $x$.
\end{definiceBox}

Symboly $-\infty,+\infty$ nejsou reálná čísla a neoznačují krajní body,
které by šly zahrnout. U nekonečna je vždy kulatá závorka.
Celé $\mathbb R$ zapíšeme $(-\infty;+\infty)$.

\begin{prikladReseniBox}[title={P2: převod oběma směry},breakable=false]
a) Zapiš $x\le-1$ intervalem a zakresli na osu.
b) Přečti $(\frac12;+\infty)$ jako nerovnost.
c) Zapiš přípustné délky součástky z úvodu.
\tcblower
a) $(-\infty;-1\rangle$; zahrnujeme $-1$ a všechna menší čísla.
\begin{center}
\begin{tikzpicture}[x=1.4cm,y=0.7cm]
\draw[->] (-4,0)--(2,0);
\foreach \x in {-3,-2,-1,0,1}{\draw (\x,-0.08)--(\x,0.08) node[below=4pt] {$\x$};}
\draw[<-,line width=1.4pt] (-4,0)--(-1,0);
\fill (-1,0) circle (2.8pt);
\end{tikzpicture}
\end{center}
b) $x>\frac12$, $x\in\mathbb R$; polovina nepatří.
c) $d\in\langle9{,}8;10{,}2\rangle$ v milimetrech.
\end{prikladReseniBox}

\begin{prikladReseniBox}[title={P3: obor nelze vynechat},breakable=false]
Je zápis $\langle-1;2\rangle$ přesným popisem \emph{celých} čísel mezi
$-1$ a $2$ včetně? Jaká celá čísla dané podmínce vyhovují?
\tcblower
Není. Interval obsahuje také například $\frac12$, které není celé.
Celá čísla splňující podmínku jsou pouze $-1,0,1,2$.
Interval vždy označuje všechna \emph{reálná} čísla mezi mezemi.
\end{prikladReseniBox}

\blok{Samostatná práce}
\begin{zadaniReseniBox}[title={S1: zápis a obrázek},breakable=false]
U každé položky zapiš druhou podobu (interval nebo nerovnost)
a nakresli osu s vyznačenými kraji. Všude je $x\in\mathbb R$.
\begin{enumerate}[label=\alph*)]
\item $-4<x\le1$;\quad b) $0\le x<\frac32$;
\item[c)] $x>-2$;\quad d) $x\le0$;
\item[e)] $(-3;-1)$;\quad f) $\langle2;5\rangle$.
\end{enumerate}
\tcblower
a) $(-4;1\rangle$: prázdný kraj $-4$, plný $1$.
b) $\langle0;\frac32)$: plný $0$, prázdný $\frac32$.
c) $(-2;+\infty)$: prázdný $-2$, doprava bez konce.
d) $(-\infty;0\rangle$: doleva bez konce, plný $0$.
e) $-3<x<-1$: oba kraje prázdné.
f) $2\le x\le5$: oba kraje plné.
\end{zadaniReseniBox}

\begin{zadaniReseniBox}[title={S2: zkontroluj význam},breakable=false]
a) Která z čísel $-2,-\frac32,0,2$ patří do $(-2;2\rangle$?
b) Oprav zápis $\langle-\infty;4\rangle$ pro $x\le4$.
c) Žák zapsal $2<x\le5$ jako $\langle2;5)$. Najdi obě chyby.
\tcblower
a) Patří $-\frac32,0,2$; $-2$ nepatří.
b) Správně $(-\infty;4\rangle$.
c) U $2$ musí být kulatá a u $5$ lomená závorka: $(2;5\rangle$.
\end{zadaniReseniBox}

\poznamka{Závěrečná kontrola: žák zdůvodní každé kolečko dosazením
krajního čísla do nerovnosti. Sjednocení ani průnik intervalů dnes nezaváděj.}

\newpage
\blok{Rezerva}
\begin{zadaniReseniBox}[title={R1: volitelně, 10–15 minut},breakable=false]
\begin{enumerate}[label=\alph*)]
\item Vypiš všechna celá čísla ležící v $(-\frac52;\frac72\rangle$.
Kolik jich je? Proč to není úplný výčet prvků intervalu?
\item Najdi racionální i iracionální číslo v $(1;\frac32)$.
\item Zapiš interval, který má levý kraj $-3$, pravý kraj $2$,
obsahuje $2$, ale neobsahuje $-3$. Má největší a nejmenší číslo?
\item Zapiš podmínkou všechna reálná čísla v $(-\infty;-0{,}5)$.
Patří tam $-\frac23$? Zdůvodni přesným porovnáním.
\end{enumerate}
\tcblower
a) $-2,-1,0,1,2,3$, celkem šest. Interval obsahuje i necelá čísla,
například $\frac14$, a má nekonečně mnoho prvků.
b) Například $\frac54$ a $\sqrt2$, protože $1<\sqrt2<1{,}5$.
c) $(-3;2\rangle$. Největší je $2$. Nejmenší neexistuje:
mezi $-3$ a kterýmkoli zvoleným prvkem leží další prvek intervalu.
d) $x<-\frac12$; ano, $-\frac23=-\frac46<-\frac36=-\frac12$.
\end{zadaniReseniBox}

\blok{Domácí procvičování}
\begin{zadaniReseniBox}[title={D1: přibližně 15–20 minut},breakable=false]
\begin{enumerate}[label=\arabic*)]
\item Zapiš intervalem a zakresli: $-5\le x<0$ a $1<x<4$, $x\in\mathbb R$.
\item Zapiš intervalem a zakresli: $x\ge\frac32$ a $x<-1$, $x\in\mathbb R$.
\item Přepiš jako nerovnosti a zakresli $(-2;6\rangle$ a $(-\infty;3)$.
\item Z čísel $-1,0,\frac12,1$ vyber ta, která patří do $\langle0;1)$.
\item Vypiš všechna celá čísla z $(-3;2)$ a vysvětli, proč jimi interval nevyčerpáš.
\item Přípustná teplota je vyšší než $-10\,^{\circ}\mathrm C$ a nejvýše
$5\,^{\circ}\mathrm C$. Zapiš rozsah intervalem a ověř obě krajní teploty.
\end{enumerate}
\tcblower
1. $\langle-5;0)$ (plný, prázdný); $(1;4)$ (oba prázdné).
2. $\langle\frac32;+\infty)$ (plný kraj, doprava);
$(-\infty;-1)$ (doleva, prázdný kraj).
3. $-2<x\le6$ (prázdný, plný); $x<3$ (doleva, prázdný kraj).
4. $0,\frac12$.
5. $-2,-1,0,1$; interval navíc obsahuje i necelá čísla, například $0{,}1$.
6. $(-10;5\rangle$ ve stupních Celsia; $-10$ nevyhovuje, $5$ vyhovuje.
\end{zadaniReseniBox}

\end{document}
